\section{Introduction}\label{sec:intro}

The idea that the Standard Model might be read off from the normed division algebras is old.  G\"unaydin and G\"ursey related the
octonions to the colour structure of quarks~\cite{gunaydin}; Dixon proposed \(\mathbb C\otimes\HH\otimes\OO\) as the algebra behind one
generation~\cite{dixon}; Furey derived the charge pattern of a generation from a number operator on \(\mathbb C\otimes\OO\)~\cite{furey2015}
and the gauge group \(SU(3)_C\times SU(2)_L\times U(1)_Y\) as a symmetry of ladder operators on \(\mathbb C\otimes\HH\otimes\OO\)~\cite{furey2018};
Gresnigt gave braided realizations~\cite{gresnigt}; and Boyle proposed triality as the origin of three generations~\cite{boyle}.  What these
constructions take as given is the role of the complex unit and the choice of chirality.  Both are usually fixed by hand.

A second, independent line starts from electromagnetic duality.  Writing the electromagnetic field as the Riemann--Silberstein vector
\(z=E+iB\)~\cite{bialynicki}, the proper Lorentz group acts on \(z\in\mathbb C^3\) as \(SO(3,\mathbb C)\), and the two field invariants
combine into one complex number \(s=z\cdot z=E^2-B^2+2iE\cdot B\).  The duality rotation \(z\mapsto e^{i\alpha}z\) turns electric into
magnetic fields and rotates \(s\); the quarter turn \(z\mapsto iz\) changes the sign of \(E^2-B^2\).  Electric- and magnetic-dominated
configurations are thereby exchanged, while null fields, \(s=0\), are fixed.

This paper connects the two lines through a single object, which we call the \emph{kernel}: a complex quaternion
\(q=it\,\mathbf 1+z_1e_1+z_2e_2+z_3e_3\in\mathbb C\otimes\HH\).  Its natural norm \(q\bar q=-t^2+s(z)\) contains exactly the Lorentz invariant
\(s\), and its unit group acts on the three complex coordinates by \(SO(3,\mathbb C)\).  We take seriously a reading in which the sign of \(s\)
separates a matter-type sector (\(s>0\)), a tachyon-type sector (\(s<0\)) and a light-like sector (\(s=0\)), and in which matter and tachyon
are the same configuration seen in two duality frames.  The working hypothesis behind the paper is that these three sectors are three
views of one object.  We do not derive physics from this hypothesis.  We ask what structure it carries, and how far that structure
reaches into the Standard Model.

\paragraph{Main results.}
\begin{enumerate}
\item In the kernel, the symmetric triple of coordinates of equal modulus at \(120^\circ\) is exactly the centred null configuration, and the
projectivized null cone is \(\mathbb CP^1\), the first base of the Hopf cascade (Section~\ref{sec:kernel}).
\item Matter-type and tachyon-type kernels of equal \(|s|\) form one orbit of \(SO(3,\mathbb C)\times U(1)_{\rm dual}\), with the light-like
kernels as the unique orbit fixed by duality.  A projector \(C=\tfrac12(1+J_B)\), with the antilinear involution \(J_Bz=i\bar z\), splits every
field into two ports; on the real doublet \((E,B)\) it becomes an orthogonal projection, and the field state becomes a point of a disk with
matter and tachyon on opposite sides and light at the centre (Section~\ref{sec:frames}).
\item The operators that commute with this involution and with \(C\) on the doublet are exactly the weightings of the two ports.  We state the
additional hypotheses under which \(C\) is the spectral projector of such an operator in general (Section~\ref{sec:cm}).
\item On \(\mathbb C\otimes\HH\otimes\OO\), a projector of the same form, \(C_{\rm ch}=\tfrac12(1+iL_{e_7})\), completes the weak doublets to the
full left-handed content of one generation with right-handed neutrino.  Among the four projectors compatible with colour and charge,
anomaly cancellation alone selects this one, up to the frame change that maps it to its mirror (Section~\ref{sec:generation}).
\item The only numerical input of the framework, a frame rapidity \(\eta\) with \(\tanh\eta=\beta_{\rm sym}\approx0.549\), can be read as the ratio
of two horizon-state countings of loop quantum gravity, selected by saturation of the uncertainty bound.  With it we state four predictions
that were registered before the relevant data: \(\sin^2\theta_{12}=0.3086\), \(m_W=80.3615\pm0.0018\) GeV, \(\sin^2\theta_{23}\le0.54\), and
\(\sin^2\theta_{23}\in[0.49,0.50]\) (Section~\ref{sec:pheno}).
\end{enumerate}

\paragraph{Limitations.}  The framework has no dynamics.  There is no action, varied in all fields, from which the kernel, the projectors or
\(\beta_{\rm sym}\) follow; the generation of item~4 is obtained under a stated setting, and the number in item~5 is an input whose counting
reading favours the less established of two treatments.  Section~\ref{sec:discussion} discusses these points and states what a dynamics would
have to provide.  An earlier, much longer version of this work contained claims that do not hold, among them steps that relied on a claimed
proof of the Riemann hypothesis; Appendix~\ref{app:earlier} lists them.

\paragraph{Outline.}  Section~\ref{sec:kernel} introduces the kernel and its Lorentz structure, Section~\ref{sec:frames} the duality frames and
the compensator, Section~\ref{sec:cm} the compensator matrix, Section~\ref{sec:generation} the generation, and Section~\ref{sec:pheno} the frame
offset and the predictions.  Section~\ref{sec:discussion} discusses limitations and open problems, and Section~\ref{sec:conclusions} concludes.
Appendix~\ref{app:method} describes how the results were checked, and Appendix~\ref{app:qam} gives the key statements in the coded notation
used in the longer version~\cite{companion}.
